\documentclass[11pt, a4paper]{article}

\usepackage[T1]{fontenc}
\usepackage[utf8]{inputenc}
\usepackage{tgpagella}
\usepackage[spanish, es-noshorthands]{babel}
\usepackage{amsmath, amssymb, bm}
\usepackage{tabularx}
\usepackage{booktabs}
\usepackage[most]{tcolorbox}
\usepackage[margin=2.5cm]{geometry}
\usepackage{ragged2e}
\usepackage{fancyhdr}

\pagestyle{fancy}
\fancyhf{}
\setlength{\headheight}{16pt}
\lhead{\textbf{UCB Sede Tarija} | Ing. Mecatrónica}
\chead{}
\rhead{IMT-121 Circuitos Electrónicos I | Corrección Autónoma}
\lfoot{Prof: M.Sc. Ing. Bernardo Quiroga Turdera}
\rfoot{Página \thepage}
\renewcommand{\headrulewidth}{0.4pt}

\definecolor{ucbblue}{RGB}{0, 51, 102}

\begin{document}

\begin{tcolorbox}[colback=ucbblue!5!white, colframe=ucbblue, title=\textbf{Actividad Paralela: Tarea de Corrección y Reflexión EC2}]
    \centering \textbf{Consolidación Individual de Competencias}
\end{tcolorbox}

\noindent\textit{Instrucciones: Mientras se desarrollan las auditorías selectivas, seleccione un problema, concepto o resultado de sus prácticas o del examen donde haya experimentado debilidad. Complete este esquema de corrección.}

\vspace{0.4cm}
\begin{enumerate}
    \item \textbf{Problema o Concepto Seleccionado:} \\
    \rule{\textwidth}{0.4pt}
    
    \item \textbf{Mi Enfoque Original (Qué hice mal):} \\
    \rule{\textwidth}{0.4pt}
    
    \item \textbf{Categoría de mi Error (Marque una):} \\
    [ ] Conceptual \quad [ ] Procedimental \quad [ ] Signos/Referencia \quad [ ] Interpretación
    
    \item \textbf{El Razonamiento Corregido (Paso a Paso):} \\
    \rule{\textwidth}{0.4pt} \\
    \rule{\textwidth}{0.4pt} \\
    \rule{\textwidth}{0.4pt}
    
    \item \textbf{Método de Verificación (Cómo me aseguro que ahora está bien):} \\
    \rule{\textwidth}{0.4pt}
    
    \item \textbf{Transferencia:} ¿En qué otro tipo de circuito o teorema podría repetirse este mismo error si no lo evito? \\
    \rule{\textwidth}{0.4pt}
\end{enumerate}

\vspace{0.5cm}
\section*{Síntesis del Flujo de Ingeniería (EC2)}
Complete la cadena lógica rellenando los espacios en blanco con la magnitud o acción principal de cada etapa:

\vspace{0.2cm}
Red Lineal $\rightarrow$ Aplicar Teoremas $\rightarrow$ Obtener \rule{3cm}{0.4pt} $\rightarrow$ Conectar \rule{3cm}{0.4pt} $\rightarrow$ Extraer \rule{3cm}{0.4pt} $\rightarrow$ Comparar Medición vs. \rule{3cm}{0.4pt}.

\end{document}
